% BOZZA della Sezione 2 (Continual Learning), versione 2. Non ancora inserita in main.tex.
% Solo CL: nessun riferimento a sicurezza o attacchi.
% Fonti verificate sui PDF in literature/papers/background/:
%   van de Ven et al. (scenari, Fig. 1, Tab. 1), Wang et al. Sez. 4 (cinque famiglie di metodi), De Lange et al. (catastrophic forgetting),
%   Wang et al. (metriche AA, FM, valutazione multi-head/single-head),
%   Lopez-Paz e Ranzato, GEM (BWT e FWT, Eq. 3-4 di LopezPaz-NeurIPS2017.pdf).
% \mycomment{...} segna cio' che manca o va verificato.
% Voci van2022three e lopez2017gem aggiunte a bibliography.bib.

\section{Continual Learning}
\label{sec:cl}

\subsection{Problem Formulation}
Continual learning (CL) addresses the problem of learning from data whose distribution changes over time, under the constraint that previously observed data cannot be accessed in their entirety~\cite{van2022three}.
In this subsection, we formalize this setting by introducing, in order, the sequence of tasks, the model, the learning objective, and the generic update rule of a CL algorithm.
This notation is adopted throughout the remainder of the paper.

\myparagraph{Tasks and data}
We model the non-stationarity of the data as a sequence of $T$ tasks $\mathcal{T}_1,\dots,\mathcal{T}_T$.
Each task $\mathcal{T}_t$ is characterized by a joint distribution $p_t(x,y)$ over the input space $\mathcal{X}$ and the label space $\mathcal{Y}$, from which its training set $\mathcal{D}_t=\{(x_i,y_i)\}_{i=1}^{N_t}$ is drawn; test samples are assumed to follow the same distribution~\cite{Wang2024tpami_cl}.
The index $t \in \{1,\dots,T\}$ denotes the task identity.
Tasks are presented sequentially: at step $t$, the learner observes only $\mathcal{D}_t$, whereas the training sets of the previous tasks are no longer available, except for the information explicitly retained by the learning algorithm.

\myparagraph{Model}
Let $f_\theta$ denote a model parametrized by a generic parameter vector $\theta$, and let $\ell$ be a loss function.
The expected loss of the model on task $\mathcal{T}_t$ is defined as a function of $\theta$:
\begin{equation}
    \mathcal{L}_t(\theta) = \mathbb{E}_{(x,y)\sim p_t}\left[\ell\left(f_\theta(x), y\right)\right].
\end{equation}
The output space of $f_\theta$ is not fixed at this stage, since it depends on the CL scenario (Sec.~\ref{sec:cl_scenarios}).

\myparagraph{Learning objective}
After step $t$, the goal is to obtain parameters that perform well on all the tasks observed so far, \ie to solve
\begin{equation}
    \min_{\theta} \sum_{j=1}^{t} \mathcal{L}_j(\theta).
\end{equation}
If $\mathcal{D}_1,\dots,\mathcal{D}_t$ were jointly available, this problem could be addressed by standard joint training.
In CL, however, the problem must be solved incrementally, using only $\mathcal{D}_t$ and the information retained from the previous steps.
The objective therefore entails a trade-off between stability, \ie preserving the performance on previous tasks (terms with $j<t$), and plasticity, \ie the ability to learn new tasks, starting from the current one (term $j=t$)~\cite{Wang2024tpami_cl}.
Minimizing $\mathcal{L}_t$ alone typically degrades the performance on previous tasks, a phenomenon known as catastrophic forgetting~\cite{de2021_cl}.

\myparagraph{Update rule}
To mitigate forgetting, a CL algorithm may retain information about previous tasks, denoted by $\mathcal{M}_{t-1}$, such as a subset of past samples or estimates of parameter importance.
A CL algorithm is then formalized as an update rule $\mathcal{A}$ that, starting from initial parameters $\theta_0$ and $\mathcal{M}_0=\emptyset$, computes at each step the parameters $\theta_t$ obtained after learning task $\mathcal{T}_t$ and the new retained information:
\begin{equation}
    (\theta_t, \mathcal{M}_t) = \mathcal{A}\left(\theta_{t-1}, \mathcal{D}_t, \mathcal{M}_{t-1}\right).
\end{equation}
The performance on task $\mathcal{T}_j$ after step $t \ge j$ is thus given by $\mathcal{L}_j(\theta_t)$.
CL methods are characterized by the content of $\mathcal{M}_{t-1}$ and by the way in which $\mathcal{A}$ exploits it (Sec.~\ref{sec:cl_methods}).

\subsection{Scenarios}
\label{sec:cl_scenarios}

CL scenarios are commonly defined for classification problems in which the tasks are non-overlapping~\cite{van2022three}.
In this setting, each sample is described by an input $x \in \mathcal{X}$, a within-task label $y \in \mathcal{Y}$, and a task label $c \in \mathcal{C}=\{1,\dots,T\}$ identifying the task to which it belongs.
Three scenarios are distinguished according to the mapping that the model must learn (Table~\ref{tab:scenarios}), which determines whether the task identity is available at test time and, if it is not, whether it must be inferred~\cite{van2022three}.

\begin{table}[t]
\centering
\small
\begin{tabular}{lccc}
\toprule
 & \textbf{Task label at test} & \textbf{Mapping learned} & \textbf{Outputs} \\
\midrule
TIL & given & $\mathcal{X}\times\mathcal{C}\rightarrow\mathcal{Y}$ & per task \\
DIL & not given, not needed & $\mathcal{X}\rightarrow\mathcal{Y}$ & shared \\
CIL & not given, inferred & $\mathcal{X}\rightarrow\mathcal{C}\times\mathcal{Y}$ & all classes \\
\bottomrule
\end{tabular}
\caption{The three CL scenarios~\cite{van2022three}.}
\label{tab:scenarios}
\end{table}

\myparagraph{Task-incremental learning (TIL)}
The task label $c$ is available at test time, so it is always clear which task must be solved.
The model can therefore use task-specific components, such as a separate output layer for each task, and the output space of a sample is limited to the classes of its own task.

\myparagraph{Domain-incremental learning (DIL)}
The task label is not available at test time, but it is not needed: every task has the same output space and only the input distribution changes across tasks, for example the same objects observed under different lighting conditions.
Task-specific components can therefore be used only if the task is first identified, which the problem does not require; consequently, unlike TIL, forgetting cannot be prevented by design~\cite{van2022three}.

\myparagraph{Class-incremental learning (CIL)}
The task label is not available at test time and must be inferred: the model must discriminate among all the classes seen so far, including classes that were never observed together.
For example, a model that first learns cats and dogs and later cows and horses must tell a cat from a cow, not only a cat from a dog.

% \myparagraph{Other settings}
% Two settings relax the task structure.
% In online CL, the training data of each task arrive as a one-pass stream.
% In task-free CL, task identities are provided neither in training nor in testing~\cite{Wang2024tpami_cl}.

\subsection{Families of CL Methods}
\label{sec:cl_methods}
CL methods differ in how they limit forgetting, that is, in what they retain about the past ($\mathcal{M}_{t-1}$) and in how the update rule $\mathcal{A}$ uses it.
Following~\cite{Wang2024tpami_cl}, we group them into five families: regularization-based, replay-based, optimization-based, representation-based, and architecture-based methods.
\mycomment{Representative methods with citations, if needed by the later sections.}

\myparagraph{Regularization-based}
These methods add an explicit regularization term to the loss to balance old and new tasks, which usually requires storing a frozen copy of the previous model~\cite{Wang2024tpami_cl}.
The term either penalizes the variation of each parameter according to its importance for the old tasks (weight regularization), or preserves the outputs of the previous model through knowledge distillation (function regularization).

\myparagraph{Replay-based}
These methods approximate and recover the data distributions of the old tasks~\cite{Wang2024tpami_cl}.
The replayed content can be a small buffer of stored training samples (experience replay), samples produced by a generative model (generative replay), or feature representations instead of raw data (feature replay).

\myparagraph{Optimization-based}
Instead of adding terms to the loss, these methods explicitly design and manipulate the optimization procedure~\cite{Wang2024tpami_cl}.
Typical examples constrain the parameter updates of the current task so that they do not interfere with the old tasks (gradient projection), or rely on meta-learning or on the geometry of the loss landscape.

\myparagraph{Representation-based}
These methods create and exploit robust representations for CL, mainly through self-supervised learning and large-scale pre-training~\cite{Wang2024tpami_cl}.
This family includes the methods that adapt a pre-trained model, kept fixed, to the downstream tasks with a small set of additional parameters, such as prompts or adapters.

\myparagraph{Architecture-based}
Whereas the previous families learn all tasks with a shared set of parameters, these methods construct task-specific parameters to avoid inter-task interference~\cite{Wang2024tpami_cl}.
The dedicated parameters can be allocated within a fixed network or obtained by expanding the architecture.

\subsection{Evaluation Metrics}
Evaluating a continual learner requires measuring more than final accuracy, because a model can score well on the last task while having forgotten the earlier ones.
CL metrics therefore cover three aspects~\cite{Wang2024tpami_cl}: the overall performance on the tasks learned so far, the stability on old tasks, and the plasticity on new tasks.
All of them are computed from the accuracies measured on each task after each learning step, which we define first.

\myparagraph{Accuracy matrix}
Let $a_{t,j} \in [0,1]$ be the accuracy on the test set of task $\mathcal{T}_j$ after learning task $\mathcal{T}_t$, for any $j \in \{1,\dots,T\}$, including tasks not yet learned~\cite{lopez2017gem}.
All the metrics below are computed from these values~\cite{Wang2024tpami_cl}.

\myparagraph{Prediction space}
The output space used to compute $a_{t,j}$ depends on the scenario.
In multi-head evaluation (TIL), the prediction is restricted to the classes of task $\mathcal{T}_j$; in single-head evaluation (CIL), it ranges over all the classes seen up to task $\mathcal{T}_t$~\cite{Wang2024tpami_cl}.

\myparagraph{Overall performance}
The average accuracy (AA) after task $\mathcal{T}_t$ is
\begin{equation}
    \mathrm{AA}_t = \frac{1}{t}\sum_{j=1}^{t} a_{t,j}.
\end{equation}

\myparagraph{Stability}
Forgetting of task $\mathcal{T}_j$ is the drop from the best accuracy it ever reached,
\begin{equation}
    f_{t,j} = \max_{i \in \{1,\dots,t-1\}} a_{i,j} - a_{t,j}, \quad j < t,
\end{equation}
and the forgetting measure (FM) averages it over the old tasks:
\begin{equation}
    \mathrm{FM}_t = \frac{1}{t-1}\sum_{j=1}^{t-1} f_{t,j}.
\end{equation}
Backward transfer (BWT)~\cite{lopez2017gem} measures the average influence of learning task $\mathcal{T}_t$ on the old tasks, and forgetting appears as a negative value:
\begin{equation}
    \mathrm{BWT}_t = \frac{1}{t-1}\sum_{j=1}^{t-1} \left(a_{t,j} - a_{j,j}\right).
\end{equation}

\myparagraph{Plasticity}
Forward transfer (FWT)~\cite{lopez2017gem} measures the influence of the previous tasks on the performance on a task not yet learned.
Let $b_j$ be the accuracy on the test set of task $\mathcal{T}_j$ at random initialization; FWT compares it with the accuracy on $\mathcal{T}_j$ just before learning it:
\begin{equation}
    \mathrm{FWT}_t = \frac{1}{t-1}\sum_{j=2}^{t} \left(a_{j-1,j} - b_j\right).
\end{equation}
